\section{Verify: 복원식과 최종 승인 predicate}
\label{sec:verify}

\paragraph{검증 당위성.}
Verify는 공격자가 선택할 수 있는 signature와 message를 받아 승인 여부를
결정하는 외부 신뢰 경계이다. 복원식의 parity, arithmetic shift, centered
reduction 또는 hash transcript가 틀리면 honest signature를 거부할 수 있고,
boolean 판정이 틀리면 명세의 norm·challenge 조건 밖 object를 승인할 수 있다.
후자가 곧 EUF-CMA 위조를 증명하는 것은 아니지만, 구현에 명세의 보안 정리를
적용할 전제를 없앤다. 그러므로 단순히 ``정상 입력을 승인한다''는 충분조건만
증명하지 않고, actual \(reject=0\)과 명세 predicate의 필요충분 동치를
증명해야 한다. parser 성공을 전제로 하는 현재 경계도 정리문에 공개하여
malformed-input 검증과 혼동하지 않는다.

\subsection{파싱 이후의 signature object}

이 장의 주 정리는 \proc{_unpack_sig_full}이 성공하여 canonical object
\((x,v,h,c)\)를 얻은 경계에서 시작한다. \(x\)는 \(z_1\)의 1024개 high
coefficient, \(v\)는 그 1024개 low byte, \(h\)는 512개 hint coefficient,
\(c\)는 weight \(58\) challenge이다.

\proc{_verify_prepare_z1_wprime}가 다음을 계산함을 증명한다.

\begin{align}
\widetilde z_1 &=256\,\Decode(x)+v, \tag{V-1}\\
w'&=\LSB(\widetilde z_0-c). \tag{V-2}
\end{align}

검증키 \((seed_A,b_1)\)에서 생성되는 네 열의 행렬을

\[
  A_1=\left(2(a-2b_1)+qj\;\middle|\;2A_{\mathrm{gen}}\right)
  \tag{V-3}
\]

로 둔다. \proc{_verify_matrix_crt}와
\proc{_sign_verify_recover_w_z2}가 다음 값을 계산함을 증명한다.

\begin{align}
u &= A_1\widetilde z_1-qcj, \tag{V-4}\\
w_1 &= h+\HighBits_{512}(u)\pmod{252}, \tag{V-5}\\
\widetilde z_2
 &=\centerq\!\left(\frac{512w_1+w'j-u}{2}\right). \tag{V-6}
\end{align}

\textup{(V-6)}의 분자는 coefficientwise 짝수여야 하며, 실제 arithmetic shift와
centered reduction이 위 식과 같다는 word/integer bridge를 포함한다.

\subsection{승인 조건}

실제 mode-2 호출은 norm bound \(163265017\)과 challenge weight \(58\)을
사용한다. \proc{_sign_verify_norm_reject}와
\proc{_sign_verify_tail_m23}에 대해 다음 동치가 목표이다.

\begin{align}
reject=0 \quad\Longleftrightarrow\quad
&\norm{\widetilde z_1}_2^2+
 \norm{\widetilde z_2}_2^2\le 163265017
\notag\\[-0.2em]
&\land\;
c=\SampleInBall\!\left(
 H(w_1,w',\widetilde\mu),58
\right),
\tag{V-7}
\end{align}

여기서

\[
  \widetilde\mu=H_{\mathrm{gen}}(vk[0..992],pre,M)
  \tag{V-8}
\]

는 실제 Verify descriptor와 memory가 지정한 byte transcript이다.

\begin{verificationgoal}[Verify core theorem]
\(siglen=1474\), mode-2 검증키, canonical parsed signature object 및 정확한
descriptor/message memory 계약 아래, 실제
\proc{_verify_prepare_z1_wprime}부터
\proc{_sign_verify_tail_m23}까지의 bounded 실행이
\textup{(V-1)}--\textup{(V-8)}을 만족하고, 반환값 \(reject=0\)이
\textup{(V-7)}의 두 수학적 조건과 동치임을 증명한다.
\end{verificationgoal}

전체 \proc{_verify_full_m23}에 대한 임의 malformed byte 입력의 동치는 별도
목표이다. 이 장은 parser 성공을 전제로 숨기지 않고 theorem precondition에
명시한다.
